Global aviation accounts for approximately 2--3\% of anthropogenic CO\textsubscript{2} emissions. At the same time, it remains one of the most challenging sectors to decarbonize due to the high energy density requirements of jet fuel and the limited feasibility of electrification for long-haul flights. Recognizing this challenge, the International Air Transport Association (IATA) has set a cap on net aviation CO\textsubscript{2} emissions starting from 2020 and aims for a 50\% reduction in emissions by 2050 relative to 2005 levels~\cite{iata2021netzero}.

Among decarbonization pathways, such as liquefied Natural Gas (NG), hydrogen, electrofuels, and electrification, Sustainable Aviation Fuels (SAF) have emerged as one of the most viable solutions due to their drop-in compatibility with existing aviation infrastructure and economic competitiveness~\cite{atsonios2015alternative}. SAF typically refers to synthesized paraffinic kerosene derived from biomass and blended with conventional jet fuel to meet performance and regulatory standards~\cite{doliente2020bio}.

The currently dominant SAF production pathways are limited by resource scarcity and harsh operating conditions. For instance, the Hydroprocessed Esters and Fatty Acids (HEFA) method dominates current SAF markets due to its simplicity and compatibility with existing refinery infrastructure. However, HEFA relies on limited feedstocks such as waste oils and animal fats, which face supply constraints as demand grows. Projections indicate that HEFA could supply only 70\% of 2030 SAF demand due to feedstock scarcity and competition with food and feed markets~\cite{doliente2020bio}. Gasification Fischer-Tropsch (FT) synthesis provides another proven pathway but is hindered by the prohibitive high temperature of gasification (1650~$^{\circ}$F) and syngas contaminants such as tars and alkali metals, which degrade catalyst performance~\cite{jones2013process}. These challenges call for innovative approaches that leverage abundant, non-food biomass while minimizing infrastructure costs.

Pyrolysis of lignocellulosic biomass offers several key advantages that address limitations in existing SAF production technologies. First, it utilizes abundant Forestry Residues (FR) and agricultural wastes, which have the potential to meet current SAF demand. These resources constitute over 60\% of global biomass potential but remain largely underutilized~\cite{wang2019biomass}. Regions such as South Korea, with forests covering 63\% of its land area, provide a sustainable feedstock supply for scaling SAF production. 
Second, pyrolysis produces high-value aromatic and cycloalkane hydrocarbon from lignin, which is typically discarded or burned in the conventional pulp industry~\cite{son2022hybrid}. Aromatic and cycloalkane hydrocarbon compounds are essential for producing jet-range fuels with superior cold flow properties and combustion stability~\cite{escalante2022evaluation}. By combining feedstock abundance and lignin valorization, pyrolysis establishes itself as a promising pathway for SAF production.

The deployment of biomass pyrolysis systems must critically evaluate the non-linear interdependency between process scale-up and logistical constraints. While increasing the capacity of integrated biorefinery theoretically yields superior economies of scale, recent techno-economic assessments demonstrate that these benefits are frequently marginalized, or even negated, by escalating primary transport costs associated with long-distance biomass haulage~\cite{chireshe2025cost}. This scale-up paradox is intensified by the heterogeneous spatial distribution and varying supply availability of feedstock. Consequently, a robust deployment strategy must balance the unit-cost reductions via economies of scale against the regional logistics of the biomass supply chain, to optimize overall cost-effectiveness in SAF production process.

Beyond logistical constraints, the chemical upgrading of pyrolysis-derived bio-oil presents significant environmental and economic hurdles. Specifically, converting oxygenated bio-oil into jet-range hydrocarbons requires hydrogen gas in an amount greater than one-third of the dry oil feed by mass~\cite{jones2013process}. Currently, this hydrogen is predominantly produced via Steam Methane Reforming (SMR), a process associated with high carbon intensity—releasing approximately 4–8 kg of CO\textsubscript{2} equivalent per kg of H\textsubscript{2}—even when coupled with Carbon Capture and Storage (CCS) technologies~\cite{IEA2023HydrogenDefinitions}. This reliance undermines the GHG reduction benefits of SAF and maintains dependence on fossil-based markets. To unlock the full decarbonization potential, alternative interventions are necessary. These include the adoption of green ammonia—a high-density carrier capable of on-demand cracking into H\textsubscript{2}~\cite{valera2018ammonia, afif2016ammonia}—or the integration of CCS units to treat post-combustion gases directly from the SAF production process.

Therefore, this paper assesses the process scale-up economics by explicitly modeling resource scarcity and its resulting price dynamics across these diverse configurations. By integrating Techno-Economic Analysis (TEA) and Life Cycle Assessment (LCA), the study determines how feedstock supply limitations govern the holistic viability of the process. The analysis specifically evaluates the economic and environmental resilience of three hydrogen sourcing pathways—conventional SMR, SMR with CCS, and green ammonia cracking—alongside varying scales of post-combustion carbon capture. Ultimately, this approach demonstrates which process designs remain robust against the backdrop of escalating feedstock prices, providing a quantitative framework for sustainable large-scale SAF deployment.

The remainder of this study is organized as follows. Section 2 details the comprehensive methodology, including the process design for pyrolysis and hydrogen pathways, the integrated TEA and LCA frameworks, and the modeling of scale-dependent feedstock price dynamics. Section 3 presents the results and discussion, evaluating the economic and environmental trade-offs of alternative hydrogen sources and carbon capture integrations under varying supply constraints. Finally, Section 4 derives the key findings to propose strategies for resilient, large-scale SAF deployment in a long-term. 

